\documentclass[10pt,a4paper]{amsart}
\usepackage[margin=19mm]{geometry}
\usepackage{amsmath,amssymb,mathtools}
\usepackage[scr=rsfs,cal=euler]{mathalfa}
\usepackage{xfrac}
\usepackage[unicode,hidelinks,bookmarks=false]{hyperref}
\newcommand{\ie}{i.e.\ }
\newcommand{\cf}{cf.\ }
\newcommand{\resp}{respectively\ }
\newcommand{\Subsection}{\subsection}
\providecommand{\nobreakdash}{\nobreak\hbox{-}}
\setlength{\emergencystretch}{2em}
\multlinegap0pt
\usepackage[shortlabels]{enumitem}
\usepackage{bm}
\usepackage{amssymb}
\usepackage{dsfont}
\usepackage{algorithm}\usepackage{algpseudocode}
\usepackage{mathtools}
\usepackage{float}
\newtheorem{assumption}{Assumption}
\newtheorem{lemma}{Lemma}
\newcommand{\Phiobj}{\Phi}
\newcommand{\R}{\mathbb{R}}
\title{BROS: Bias-Corrected Randomized Subspaces for Memory-Efficient Single-Loop Bilevel Optimization}
\author{Hengrui Zhang, Boao Kong, Engao Zhang, Kun Yuan}
\date{Source: arXiv:2605.10288v2 (CC BY 4.0)}
\begin{document}
\maketitle

\noindent\emph{Source and licence.} BROS: Bias-Corrected Randomized Subspaces for Memory-Efficient Single-Loop Bilevel Optimization. Version arXiv:2605.10288v2 (CC BY 4.0), distributed by arXiv under the Creative Commons Attribution 4.0 International licence, http://creativecommons.org/licenses/by/4.0/, as recorded in the arXiv licence metadata for this exact version and shown on the abstract page https://arxiv.org/abs/2605.10288v2. Full licence notice: \texttt{licenses/arxiv-2605.10288-CC-BY-4.0.txt}.\par\smallskip

\noindent\textit{Editorial scope.} Counted unit: the article's lemma lem:smoothness\_consequences (source0.tex lines 884--937). The article's complete original proof of it is delivered with it (source0.tex lines 939--1082, one \texttt{proof} environment of 144 lines). No mathematics is written, repaired or re-derived: the statement and the proof are the article's own lines, in the article's own notation, and no retained line is substituted or re-flowed.  Retained uncounted as the source-local context the proof needs: lines 483--492.  Nothing was omitted from any retained span, so the package carries no omission record.  No retained line contains a citation, so the wrapper carries no bibliography and no bibliography entry.  No retained line contains a figure environment, an image-inclusion command or drawing code, so no image is delivered and the package declares no asset dependency.  Editorial formatting only, in addition: an AMS \texttt{amsart} preamble that restates the article's own declarations and macros used by the retained lines, namely the environments \texttt{assumption}, \texttt{lemma} on separate counters, together with the standard packages those lines need.  The retained environments print in this document as Assumption 1, Lemma 1 in order of appearance, whereas the article prints the same items with the numbering of its own full document.  The complete native source of the article as distributed in the arXiv e-print for this exact version is bundled unchanged as \texttt{sources/200286/source0.tex}.
\begin{assumption}[Smoothness, boundedness, and lower-level curvature]
\label{as:A1}
There exist positive constants $\mu_g,L_{f,0},L_{f,1},L_{g,1},L_{g,2}>0$ such that:
\begin{enumerate}[leftmargin=1.5em,label=\arabic*.,topsep=1pt,itemsep=1pt,parsep=0pt,partopsep=0pt]
    \item $\nabla f$ and $\nabla g$ are Lipschitz continuous in $(x,\boldsymbol Y)$ with constants $L_{f,1}$ and $L_{g,1}$, and $\nabla^2 g$ is $L_{g,2}$-Lipschitz continuous as a linear operator on the product space $\mathbb R^{d_x}\times\mathcal Y$;
    \item $\|\nabla_{\boldsymbol Y} f(x,\boldsymbol Y^*(x))\|\le L_{f,0}$ for all $x$;
    \item for every fixed $x$, $g(x,\cdot)$ is $\mu_g$-strongly convex on $\mathcal Y$.
\end{enumerate}
Let $L_{\max}:=\max\{L_{f,0},L_{f,1},L_{g,1},L_{g,2}\}$ and $\kappa:=L_{\max}/\mu_g$.
\end{assumption}

\begin{lemma}[Product-space smoothness consequences]\label{lem:smoothness_consequences}
Suppose that \emph{(A1)} holds. Let $\boldsymbol Y^*(x)$ be the unique minimizer of $\boldsymbol Y\mapsto g(x,\boldsymbol Y)$ and define
\begin{equation}\label{eq:def_Zstar}
\boldsymbol Z^*(x):=\mathcal H(x,\boldsymbol Y^*(x))^{-1}[\nabla_{\boldsymbol Y}f(x,\boldsymbol Y^*(x))]\in\mathcal Y.
\end{equation}
Then the following statements hold.
\begin{enumerate}
\item \textbf{Implicit hypergradient formula.}
The bilevel objective $\Phiobj(x)=f(x,\boldsymbol Y^*(x))$ is differentiable and
\begin{equation}\label{eq:hypergrad_formula}
\nabla\Phiobj(x)=\nabla_x f(x,\boldsymbol Y^*(x))-\mathcal J(x,\boldsymbol Y^*(x))[\boldsymbol Z^*(x)].
\end{equation}

\item \textbf{Lipschitz continuity of $\boldsymbol Y^*(\cdot)$.}
The mapping $\boldsymbol Y^*(\cdot)$ is Lipschitz continuous with
\begin{equation}\label{eq:LY_star}
\|\boldsymbol Y^*(x)-\boldsymbol Y^*(\tilde x)\|\le L_{Y^*}\|x-\tilde x\|,
\quad
L_{Y^*}:=\frac{L_{g,1}}{\mu_g}\le \kappa.
\end{equation}

\item \textbf{Lipschitz continuity of $\boldsymbol Z^*(\cdot)$.}
The mapping $\boldsymbol Z^*(\cdot)$ is Lipschitz continuous with
\begin{equation}\label{eq:LZ_star}
\|\boldsymbol Z^*(x)-\boldsymbol Z^*(\tilde x)\|
\le L_{Z^*}\|x-\tilde x\|,
\quad
L_{Z^*}:=\sqrt{1+L_{Y^*}^2}\left(\frac{L_{f,1}}{\mu_g}+\frac{L_{f,0}L_{g,2}}{\mu_g^2}\right).
\end{equation}
In particular, if $\kappa\ge 1$ then $L_{Z^*}\le 2\sqrt{2}\,\kappa^3$.

\item \textbf{Lipschitz continuity of $\nabla\Phiobj$.}
The gradient $\nabla\Phiobj$ is Lipschitz continuous with
\begin{equation}\label{eq:Lnabla_Phi}
\|\nabla\Phiobj(x)-\nabla\Phiobj(\tilde x)\|
\le L_{\nabla\Phi}\|x-\tilde x\|,
\end{equation}
where one may take
\begin{equation}\label{eq:Lnabla_Phi_expr}
L_{\nabla\Phi}:=
\sqrt{1+L_{Y^*}^2}\left(
L_{f,1}+\frac{L_{f,0}L_{g,2}}{\mu_g}
+\frac{L_{g,1}L_{f,1}}{\mu_g}
+\frac{L_{g,1}L_{f,0}L_{g,2}}{\mu_g^2}
\right).
\end{equation}

\item \textbf{Uniform bound on $\|\boldsymbol Z^*(x)\|$.}
For all $x$,
\begin{equation}\label{eq:Zstar_bound}
\|\boldsymbol Z^*(x)\|\le \frac{L_{f,0}}{\mu_g}.
\end{equation}
\end{enumerate}
\end{lemma}

\begin{proof}
We endow $\R^{d_x}\times\mathcal Y$ with the product norm $\|(x,\boldsymbol Y)\|:=\sqrt{\|x\|^2+\|\boldsymbol Y\|^2}$.

\textbf{Step 1: implicit differentiation and \eqref{eq:hypergrad_formula}. }
Fix $x$ and write $\boldsymbol Y^*:=\boldsymbol Y^*(x)$. Define
\begin{equation*}
F(x,\boldsymbol Y):=\nabla_{\boldsymbol Y}g(x,\boldsymbol Y)\in\mathcal Y.
\end{equation*}
By Assumption~\ref{as:A1}, $g$ is twice continuously differentiable, so $F$ is continuously differentiable with
\begin{equation*}
\nabla_{\boldsymbol Y}F(x,\boldsymbol Y)=\mathcal H(x,\boldsymbol Y), \quad \nabla_x F(x,\boldsymbol Y)=\nabla_{\boldsymbol Y x}^2 g(x,\boldsymbol Y).
\end{equation*}
Since $\boldsymbol Y\mapsto g(x,\boldsymbol Y)$ is $\mu_g$-strongly convex in the product Frobenius norm, its Hessian is uniformly positive definite:
\begin{equation}\label{eq:H_sc_product}
\langle \boldsymbol W,\mathcal H(x,\boldsymbol Y)[\boldsymbol W]\rangle
\ge \mu_g \|\boldsymbol W\|^2,
\quad
\forall \boldsymbol W\in\mathcal Y,\ \forall (x,\boldsymbol Y).
\end{equation}
Hence $\mathcal H(x,\boldsymbol Y^*)$ is invertible. Applying the implicit function theorem to $F(x,\boldsymbol Y)=0$ at $(x,\boldsymbol Y^*)$ shows that $\boldsymbol Y^*(\cdot)$ is continuously differentiable locally, and differentiating $F(x,\boldsymbol Y^*(x))\equiv 0$ in direction $u\in\R^{d_x}$ gives
\begin{equation*}
0=\nabla_{\boldsymbol Y x}^2 g(x,\boldsymbol Y^*)[u] +\mathcal H(x,\boldsymbol Y^*)[D\boldsymbol Y^*(x)[u]],
\end{equation*}
so
\begin{equation}\label{eq:DYstar_formula_product}
D\boldsymbol Y^*(x)[u]
=-\mathcal H(x,\boldsymbol Y^*)^{-1}[\nabla_{\boldsymbol Y x}^2 g(x,\boldsymbol Y^*)[u]].
\end{equation}

Now compute the directional derivative of $\Phiobj(x)=f(x,\boldsymbol Y^*(x))$:
\begin{equation*}
D\Phiobj(x)[u] =\langle \nabla_x f(x,\boldsymbol Y^*),u\rangle +\langle \nabla_{\boldsymbol Y}f(x,\boldsymbol Y^*),D\boldsymbol Y^*(x)[u]\rangle.
\end{equation*}
Substituting \eqref{eq:DYstar_formula_product} and using that $\mathcal H(x,\boldsymbol Y^*)$ is self-adjoint on $\mathcal Y$ (hence so is its inverse) yields
\begin{equation*}
D\Phiobj(x)[u]
={}\langle \nabla_x f(x,\boldsymbol Y^*),u\rangle -\Bigl\langle
\mathcal H(x,\boldsymbol Y^*)^{-1}[\nabla_{\boldsymbol Y}f(x,\boldsymbol Y^*)],
\nabla_{\boldsymbol Y x}^2 g(x,\boldsymbol Y^*)[u]
\Bigr\rangle.
\end{equation*}
We define $\mathcal J(x,\boldsymbol Y):\mathcal Y\to\R^{d_x}$ as the adjoint of $u\mapsto \nabla_{\boldsymbol Y x}^2 g(x,\boldsymbol Y)[u]$, i.e., for all $u\in\R^{d_x}$ and $\boldsymbol W\in\mathcal Y$,
\begin{equation}\label{eq:adjoint_J_product}
\langle \nabla_{\boldsymbol Y x}^2 g(x,\boldsymbol Y)[u],\boldsymbol W\rangle
=\langle u,\mathcal J(x,\boldsymbol Y)[\boldsymbol W]\rangle.
\end{equation}
Recalling the definition \eqref{eq:def_Zstar}, we obtain
\begin{equation*}
D\Phiobj(x)[u] =\langle \nabla_x f(x,\boldsymbol Y^*)-\mathcal J(x,\boldsymbol Y^*)[\boldsymbol Z^*(x)],u\rangle,
\end{equation*}
which proves \eqref{eq:hypergrad_formula}.

\textbf{Step 2: Lipschitz continuity of $\boldsymbol Y^*(\cdot)$. }
Fix $x,\tilde x\in\R^{d_x}$. Using the optimality conditions
\begin{equation*}
\nabla_{\boldsymbol Y}g(x,\boldsymbol Y^*(x))=0, \quad \nabla_{\boldsymbol Y}g(\tilde x,\boldsymbol Y^*(\tilde x))=0,
\end{equation*}
and the strong monotonicity of $\boldsymbol Y\mapsto \nabla_{\boldsymbol Y}g(x,\boldsymbol Y)$, we get
\begin{equation*}
\begin{aligned}
\mu_g\|\boldsymbol Y^*(x)-\boldsymbol Y^*(\tilde x)\|
\le{}& \|\nabla_{\boldsymbol Y}g(x,\boldsymbol Y^*(\tilde x))-\nabla_{\boldsymbol Y}g(x,\boldsymbol Y^*(x))\| \\
={}& \|\nabla_{\boldsymbol Y}g(\tilde x,\boldsymbol Y^*(\tilde x))-\nabla_{\boldsymbol Y}g(x,\boldsymbol Y^*(\tilde x))\|.
\end{aligned}
\end{equation*}
By the $L_{g,1}$-Lipschitz continuity of $\nabla g$ in Assumption~\ref{as:A1}, the last term is at most $L_{g,1}\|x-\tilde x\|$, which proves \eqref{eq:LY_star}.

\textbf{Step 3: Lipschitz continuity and the bound of $\boldsymbol Z^*(\cdot)$. }
Define
\begin{equation*}
A(x):=\mathcal H(x,\boldsymbol Y^*(x)), \quad b(x):=\nabla_{\boldsymbol Y}f(x,\boldsymbol Y^*(x)), \quad \boldsymbol Z^*(x)=A(x)^{-1}[b(x)].
\end{equation*}
From \eqref{eq:H_sc_product}, $\|A(x)^{-1}\|\le 1/\mu_g$. Assumption~\ref{as:A1} gives $\|b(x)\|\le L_{f,0}$, so $\|\boldsymbol Z^*(x)\|\le \frac{L_{f,0}}{\mu_g}$.

Next, by Assumption~\ref{as:A1} and \eqref{eq:LY_star},
\begin{equation}\label{eq:b_Lipschitz_product}
\|b(x)-b(\tilde x)\|
\le
L_{f,1}\|(x,\boldsymbol Y^*(x))-(\tilde x,\boldsymbol Y^*(\tilde x))\|
\le
L_{f,1}\sqrt{1+L_{Y^*}^2}\|x-\tilde x\|.
\end{equation}
Similarly, since $\nabla^2 g$ is $L_{g,2}$-Lipschitz,
\begin{equation}\label{eq:A_Lipschitz_product}
\|A(x)-A(\tilde x)\|
\le
L_{g,2}\|(x,\boldsymbol Y^*(x))-(\tilde x,\boldsymbol Y^*(\tilde x))\|
\le
L_{g,2}\sqrt{1+L_{Y^*}^2}\|x-\tilde x\|.
\end{equation}
Using the resolvent identity $A(x)^{-1}-A(\tilde x)^{-1} = A(x)^{-1}(A(\tilde x)-A(x))A(\tilde x)^{-1}$, we have
\begin{equation}\label{eq:Ainv_Lipschitz_product}
\|A(x)^{-1}-A(\tilde x)^{-1}\|
\le
\frac{1}{\mu_g^2}\|A(x)-A(\tilde x)\|.
\end{equation}
Decomposing
\begin{equation*}
\boldsymbol Z^*(x)-\boldsymbol Z^*(\tilde x) = A(x)^{-1}[b(x)-b(\tilde x)] + (A(x)^{-1}-A(\tilde x)^{-1})[b(\tilde x)]
\end{equation*}
and combining \eqref{eq:b_Lipschitz_product}--\eqref{eq:Ainv_Lipschitz_product} with $\|b(\tilde x)\|\le L_{f,0}$ yields \eqref{eq:LZ_star}.

\textbf{Step 4: Lipschitz continuity of $\nabla\Phiobj$. }
From \eqref{eq:hypergrad_formula},
\begin{equation*}
\nabla\Phiobj(x) = \nabla_x f(x,\boldsymbol Y^*(x)) -\mathcal J(x,\boldsymbol Y^*(x))[\boldsymbol Z^*(x)].
\end{equation*}
Because $\nabla g$ is $L_{g,1}$-Lipschitz, the full Hessian of $g$ has operator norm at most $L_{g,1}$; in particular,
\begin{equation}\label{eq:J_bound}
\|\mathcal J(x,\boldsymbol Y)\|\le L_{g,1},
\quad
\forall (x,\boldsymbol Y)\in\R^{d_x}\times\mathcal Y.
\end{equation}
Likewise, the $L_{g,2}$-Lipschitz continuity of $\nabla^2 g$ implies
\begin{equation}\label{eq:J_Lipschitz}
\|\mathcal J(x,\boldsymbol Y)-\mathcal J(\tilde x,\tilde{\boldsymbol Y})\|
\le
L_{g,2}\|(x,\boldsymbol Y)-(\tilde x,\tilde{\boldsymbol Y})\|.
\end{equation}
Therefore,
\begin{equation*}
\begin{aligned}
\|\nabla\Phiobj(x)-\nabla\Phiobj(\tilde x)\|
\le{}&
\|\nabla_x f(x,\boldsymbol Y^*(x))-\nabla_x f(\tilde x,\boldsymbol Y^*(\tilde x))\| \\
&+
\|\mathcal J(x,\boldsymbol Y^*(x))[\boldsymbol Z^*(x)]
-\mathcal J(\tilde x,\boldsymbol Y^*(\tilde x))[\boldsymbol Z^*(\tilde x)]\|.
\end{aligned}
\end{equation*}
The first term is bounded by $L_{f,1}\sqrt{1+L_{Y^*}^2}\|x-\tilde x\|$. For the second term, add and subtract $\mathcal J(\tilde x,\boldsymbol Y^*(\tilde x))[\boldsymbol Z^*(x)]$:
\begin{equation*}
\begin{aligned}
\|\mathcal J(x,\boldsymbol Y^*(x))[\boldsymbol Z^*(x)]
-\mathcal J(\tilde x,\boldsymbol Y^*(\tilde x))[\boldsymbol Z^*(\tilde x)]\| 
\le{}&
\|\mathcal J(x,\boldsymbol Y^*(x))-\mathcal J(\tilde x,\boldsymbol Y^*(\tilde x))\|\,
\|\boldsymbol Z^*(x)\| \\
&+\|\mathcal J(\tilde x,\boldsymbol Y^*(\tilde x))\|\,
\|\boldsymbol Z^*(x)-\boldsymbol Z^*(\tilde x)\|.
\end{aligned}
\end{equation*}
Using \eqref{eq:J_Lipschitz}, \eqref{eq:J_bound}, and \eqref{eq:LZ_star} yields \eqref{eq:Lnabla_Phi} with the constant \eqref{eq:Lnabla_Phi_expr}.
\end{proof}

\end{document}
